% Part: computability
% Chapter: recursive-functions
% Section: general-recursive-functions

\documentclass[../../../include/open-logic-section]{subfiles}

\begin{document}

\olfileid{cmp}{rec}{gen}
\olsection{सामान्य पुनरावर्ती फलने}

सर्वत्र परिभाषित फलनांचा एक संच मिळवण्याचा आणखी एक मार्ग आहे.
सर्वत्र परिभाषित $f(x,\vec z)$ हे फलन \emph{नियमित} असे म्हणू, जर
नैसर्गिक संख्यांच्या प्रत्येक क्रमिका~$\vec z$ साठी असा $x$ असेल की
$f(x,\vec z) = 0$. दुसऱ्या शब्दांत, ज्या फलनांना अपरिबद्ध शोध लागू
केल्यावर सर्वत्र परिभाषित फलन मिळते, तीच नियमित फलने आहेत.
सावधगिरीने अपरिबद्ध शोध नियमित फलनांपुरताच मर्यादित करता येतो:

\begin{defn}
\ollabel{defn:general-recursive}
\emph{सामान्य पुनरावर्ती फलनांचा} संच हा विविध चलसंख्यांची नैसर्गिक
संख्यांपासून नैसर्गिक संख्यांकडे जाणारी फलने असलेला सर्वांत लहान संच
आहे: त्यात शून्य, उत्तरवर्ती आणि प्रक्षेपणे आहेत, आणि तो संयोजन, आदिम
पुनरावर्तन आणि \emph{नियमित} फलनांना लागू केलेला अपरिबद्ध शोध यांखाली
बंद आहे.
\end{defn}

प्रत्येक सामान्य पुनरावर्ती फलन सर्वत्र परिभाषित आहे, हे स्पष्ट आहे.
\olref{defn:general-recursive} आणि \olref[par]{defn:recursive-fn}
यांतील फरक असा की दुसऱ्या व्याख्येत मधल्या टप्प्यांवर आंशिक पुनरावर्ती
फलने वापरण्यास परवानगी आहे; शेवटी मिळणारे फलन सर्वत्र परिभाषित असावे,
एवढीच अट आहे. म्हणून ``सामान्य'' हा ऐतिहासिक अवशेष असलेला शब्द
दिशाभूल करणारा आहे: वरवर पाहता \olref{defn:general-recursive} ही
व्याख्या \olref[par]{defn:recursive-fn} पेक्षा \emph{कमी} सामान्य
आहे. पण सुदैवाने हा फरक केवळ भासमान आहे. व्याख्या वेगळ्या असल्या,
तरी सामान्य पुनरावर्ती फलनांचा संच आणि पुनरावर्ती फलनांचा संच एकच आहे.

\end{document}

